% 原书PDF第29页7.1.4.5.3至第33页7.1.4.5.4标题前。
\bookdeepheading{paragraph}{7.1.4.5.3}{各向同性高斯参考测度与迹方法}
假设 $X$、$Y$ 分别为 $n$ 维和 $m$ 维多元正态向量，均值为零，协方差矩阵分别为 $\Sigma_X$ 和 $\Sigma_Y$（Janzing et al. 2010）。进一步假定 $Y$ 是 $X$ 的线性变换：
\begin{equation}
Y=AX,
\tag{7.111}
\end{equation}
其中 $A$ 为 $m\times n$ 矩阵。归一化迹（renormalized trace）定义为
\begin{equation}
\tau_n(\cdot)=\frac{\Tr(\cdot)}n.
\tag{7.112}
\end{equation}
令 $O(n)$ 表示 $n\times n$ 实正交矩阵组成的群。附录7.E证明迹方法的以下基本结果（Janzing et al. 2010，定理1）。

\par\medskip\Needspace{6\baselineskip}\noindent\textbf{结论7.15：迹的乘法性}\par
设 $\Sigma$ 为 $n\times n$ 对称正定矩阵，$A$ 为 $m\times n$ 矩阵。令 $U$ 为按照Haar测度，从 $n\times n$ 实正交矩阵群 $O(n)$ 中随机选取的正交矩阵。则
\begin{equation}
\left|\tau_m(AU\Sigma U^TA^T)-\tau_n(\Sigma)\tau_m(AA^T)\right|
\le2\varepsilon\lVert\Sigma\rVert\lVert AA^T\rVert
\tag{7.113}
\end{equation}
以至少 $q=1-\exp(-\kappa(n-1)\varepsilon^2)$ 的概率成立。其中常数 $\kappa$ 独立于 $\Sigma,A,n,m,\varepsilon$，$\lVert\cdot\rVert$ 表示矩阵范数。结论7.15蕴含，矩阵对 $(A,\Sigma_X)$ 满足
\begin{equation}
\tau_n(A\Sigma_XA^T)\approx\tau_n(\Sigma)\tau_m(AA^T).
\tag{7.114}
\end{equation}
式 (7.99) 可用于推断因果关系。定义
\begin{equation}
\Delta_{X\to Y}=\log\tau_n(A\Sigma A^T)-\log\tau_n(\Sigma)-\log\tau_m(AA^T).
\tag{7.115}
\end{equation}
习题6表明，一维情形下 $\Delta$ 等于零，意味着迹方法不能用于两个一维变量。

下面证明：若假设 $n\le m$，且 $A$ 的秩为 $n$，则 $\Delta_{X\to Y}=0$，$\Delta_{Y\to X}\le0$。考虑确定性线性模型
\[
Y=AX,\qquad X=A^-Y,
\]
其中 $A^-$ 表示矩阵 $A$ 的伪逆。定义
\begin{equation}
\Delta_{Y\to X}=\log\tau_n(A^-\Sigma_{YY}A^{-T})
-\log\tau_m(\Sigma_{YY})-\log\tau_n(A^-A^{-T}).
\tag{7.116}
\end{equation}
假设实值随机变量 $Z$ 服从 $AA^T$ 的特征值经验分布，即对所有 $k\in\mathbb Z$，均有 $\tau_m((AA^T)^k)=E(Z^k)$。可以证明（习题7），
\begin{equation}
\Delta_{X\to Y}+\Delta_{Y\to X}
=-\log(1-\operatorname{cov}(Z,1/Z))+\log\frac nm.
\tag{7.117}
\end{equation}
粗略而言，由于 $E[Z]\ge0$、$E[1/Z]\ge0$、$E[Z]E[1/Z]\ge1$，可得
\[
\Cov(Z,1/Z)\le0,\qquad-\log(1-\Cov(Z,1/Z))\le0.
\]
当 $n=m$ 时，
\begin{equation}
\Delta_{X\to Y}+\Delta_{Y\to X}\le0.
\tag{7.118}
\end{equation}

下面以各向同性高斯分布作为参考测度，考虑原因 $X$ 的协方差矩阵的各向异性。先定义协方差矩阵的各向异性（Janzing et al. 2010）。

\begin{bookdefinition}{7.9}
\textbf{协方差矩阵的各向异性。}协方差矩阵 $\Sigma_1$ 的各向异性，定义为高斯分布 $\Sigma_1$ 与协方差 $\Sigma_0=\lambda I$ 的各向同性高斯分布之间的最小K–L距离：
\begin{equation}
D(\Sigma_1)=\min_\lambda D(P_{\Sigma_1}\Vert P_{\Sigma_0}).
\tag{7.119}
\end{equation}
\end{bookdefinition}
附录7.F证明，
\begin{equation}
D(\Sigma_1)=\frac12\bigl(n\log\tau_n(\Sigma_1)-\log|\Sigma_1|\bigr).
\tag{7.120}
\end{equation}
假设 $\varepsilon_X$、$\varepsilon_Y$ 都是各向同性高斯分布的流形，$\varepsilon_X^*$、$\varepsilon_Y^*$ 分别是 $P_X$、$P_Y$ 的投影。附录7.F还证明结论7.16（Janzing et al. 2010）。

\par\medskip\Needspace{5\baselineskip}\noindent\textbf{结论7.16：输出协方差矩阵的各向异性与K–L距离}\par
假设 $P_X$ 与 $P_Y$ 的协方差矩阵分别为 $\Sigma_X$ 与 $\Sigma_Y=A\Sigma_XA^T$，则
\begin{equation}
D(\Sigma_X)=D(P_X\Vert\varepsilon_X^*)
=\frac12\bigl(n\log\tau_n(\Sigma_X)-\log|\Sigma_X|\bigr),
\tag{7.121}
\end{equation}
以及
\begin{equation}
\begin{aligned}
D(A\Sigma_XA^T)&=D(P_Y\Vert\varepsilon_Y^*)
=\frac12\bigl(n\log\tau_n(\Sigma_Y)-\log|\Sigma_Y|\bigr)\\
&=\frac n2\Delta_{X\to Y}+D(\Sigma_X)+D(AA^T).
\end{aligned}
\tag{7.122}
\end{equation}
附录7.F还证明关于独立性的结论7.17（Janzing et al. 2012；Zscheischler et al. 2011）。

\par\medskip\Needspace{4\baselineskip}\noindent\textbf{结论7.17：原因与线性变换矩阵的独立性}\par
若 $X\to Y$ 且 $Y=AX$，则
\begin{equation}
\tau_n(A\Sigma_XA^T)=\tau_n(\Sigma_X)\tau_n(AA^T).
\tag{7.123}
\end{equation}

现在将 $X$、$Y$ 之间的确定性关系扩展为随机关系。考虑一般线性模型（Janzing et al. 2010）：
\begin{equation}
Y=AX+e,
\tag{7.124}
\end{equation}
其中 $A$ 为 $m\times n$ 矩阵，$e$ 为均值零、协方差矩阵 $\Sigma_e$ 的噪声向量，在统计上独立于 $X$。容易证明，
\begin{equation}
\Sigma_Y=A\Sigma_XA^T+\Sigma_e,
\tag{7.125}
\end{equation}
\begin{equation}
\widehat A=\Sigma_{YX}\Sigma_X^{-1}.
\tag{7.126}
\end{equation}
反向模型为
\begin{equation}
X=\widetilde AY+\widetilde e,
\tag{7.127}
\end{equation}
其中
\begin{equation}
\widetilde A=\Sigma_{XY}\Sigma_Y^{-1}.
\tag{7.128}
\end{equation}
若假设 $A$ 为正交变换，$e$ 为各向同性，即 $\Sigma_e=\lambda I$，则可证明结论7.18（Janzing et al. 2010；附录7.G）。

\par\medskip\Needspace{5\baselineskip}\noindent\textbf{结论7.18：含噪声线性模型的迹方法}\par
考虑 $Y=AX+e$，其中 $A$ 为正交矩阵，噪声向量 $e$ 的协方差为 $\Sigma_e=\lambda I$，$\lambda>0$。则
\begin{equation}
\begin{aligned}
\Delta_{Y\to X}
&=\log\frac{\tau_n(\widetilde A\Sigma_Y\widetilde A^T)}
{\tau_n(\Sigma_Y)\tau_n(\widetilde A\widetilde A^T)}\\
&=\log\frac1n\sum_{i=1}^n\frac{\mu_i^2}{\mu_i+\lambda}
-\log\frac1n\sum_{i=1}^n(\mu_i+\lambda)\\
&\hspace{2em}-\log\frac1n\sum_{i=1}^n\frac{\mu_i^2}{(\mu_i+\lambda)^2}>0,
\end{aligned}
\tag{7.129}
\end{equation}
其中 $\mu_1\ge\mu_2\ge\cdots\ge\mu_n\ge0$ 为协方差矩阵 $\Sigma_X$ 的特征值，且 $n>1$。

若假设 $X$ 与 $Y$ 之间的关联可由
\[
r=\frac1n\sum_{i=1}^n\lambda_i^2=\frac1n\Tr(R^2)
\]
度量，则在同一模型下，可证明刻画关联的结论7.19（附录7.H）。显然，同一模型下有
\[
\begin{aligned}
\Delta_{X\to Y}&=\log\tau_n(A\Sigma_XA^T)-\log\tau_n(\Sigma_X)-\log\tau_n(AA^T)\\
&=\log\tau_n(\Sigma_X)-\log\tau_n(\Sigma_X)-\log(1)=0.
\end{aligned}
\]

\par\medskip\Needspace{5\baselineskip}\noindent\textbf{结论7.19：关联的刻画}\par
假设模型为
\[
Y=AX+e,
\]
其中 $A$ 为可逆矩阵，噪声向量 $e$ 满足 $\Sigma_e=\lambda I$，$\lambda>0$。则 $X$ 与 $Y$ 之间的关联度量为
\begin{equation}
r=\frac1n\sum_{i=1}^n\frac{\mu_i}{\mu_i+\lambda},
\tag{7.130}
\end{equation}
且 $r<1$、$\log r<0$。

不假设 $A$ 正交时，因果与关联度量计算为
\begin{equation}
\begin{aligned}
\Delta_{X\to Y}
&=\log\tau_m(\Sigma_{YX}\Sigma_X^{-1}\Sigma_{XY})-\log\tau_n(\Sigma_X)\\
&\hspace{2em}-\log\tau_m(\Sigma_{YX}\Sigma_X^{-2}\Sigma_{XY}),
\end{aligned}
\tag{7.131}
\end{equation}
\begin{equation}
\begin{aligned}
\Delta_{Y\to X}
&=\log\tau_n\bigl(\Sigma_{XY}(\Sigma_{YX}\Sigma_X^{-1}\Sigma_{XY}+\lambda I)^{-1}\Sigma_{YX}\bigr)\\
&\hspace{2em}-\log\tau_m(\Sigma_{YX}\Sigma_Y^{-1}\Sigma_{XY}+\lambda I)\\
&\hspace{2em}-\log(\tau_n(\Sigma_{XY}(\Sigma_{YX}\Sigma_X^{-1}\Sigma_{XY}+\lambda I)^{-2}\Sigma_{YX},
\end{aligned}
\tag{7.132}
\end{equation}
以及
\begin{equation}
r=\log\tau_n(\Sigma_{YX}\Sigma_X^{-1}\Sigma_{XY}\Sigma_Y^{-1}).
\tag{7.133}
\end{equation}
综上，若 $X\to Y$，则 $\Delta_{X\to Y}=0$、$\Delta_{Y\to X}>0$；若 $Y\to X$，则 $\Delta_{X\to Y}>0$、$\Delta_{Y\to X}=0$；若 $X$ 与 $Y$ 关联，则 $r<0$。线性因果模型检验的统计过程，可汇总为结论7.20（Janzing et al. 2010）。

\par\medskip\Needspace{5\baselineskip}\noindent\textbf{结论7.20：线性因果检验的迹算法}\par
迹算法如下。

\noindent\textbf{步骤1：}估计 $\Sigma_X$、$\Sigma_Y$、$\Sigma_{XY}$ 和 $\Sigma_{YX}$。

\noindent\textbf{步骤2：}计算 $A=\Sigma_{YX}\Sigma_{XX}^{-1}$。

\noindent\textbf{步骤3：}计算 $\widetilde A=\Sigma_{XY}\Sigma_{YY}^{-1}$。

\noindent\textbf{步骤4：}计算
\[
\Delta_{X\to Y}=\log\frac{\tau_m(A\Sigma_XA^T)}{\tau_n(\Sigma_X)\tau_m(AA^T)}
\]
与
\[
\Delta_{Y\to X}=\log\frac{\tau_n(\widetilde A\Sigma_Y\widetilde A^T)}
{\tau_m(\Sigma_Y)\tau_n(\widetilde A\widetilde A^T)}.
\]
\noindent\textbf{步骤5：}若 $|\Delta_{Y\to X}|>(1+\varepsilon)|\Delta_{X\to Y}|$，则
\[
X\to Y.
\]
否则：
\begin{quote}
若 $|\Delta_{X\to Y}|>(1+\varepsilon)|\Delta_{Y\to X}|$，则
\[
Y\to X.
\]
否则：没有因果关系，然后检验关联。

结束条件判断。
\end{quote}
结束条件判断。
